\begin{figure}[p]
    \centering

    \begin{subfigure}[b]{0.48\textwidth}
        \centering
        \begin{tikzpicture}[ 
            scale=1, auto, node distance=2cm, 
            every edge/.style={draw, <-, font=\small},
            state/.style={circle, fill=mylightgrey, draw=black, text=black}
            ]
        
            % Nodes
            \node[state] (s1) at (0,2) {$s_1$};
            \node[state] (s2) at (2,2) {$s_2$};
            \node[state] (s3) at (2,0) {$s_3$};
            \node[state] (s4) at (0,0) {$s_4$};
            
            % Self-loops
            \path[->] (s1) edge[loop above, out=115, in=155, looseness=8] (s1);
            \path[->] (s2) edge[loop above, out=25, in=65, looseness=8] (s2);
            \path[->] (s3) edge[loop below, out=295, in=335, looseness=8] (s3);
            \path[->] (s4) edge[loop below, out=205, in=245, looseness=8] (s4);
            
            % Arcs between nodes
            \path[->] (s1) edge[bend left=12] (s2);
            \path[->] (s1) edge[bend left=12] (s3);
            \path[->] (s1) edge[bend left=12] (s4);
            
            \path[->] (s2) edge[bend left=12] (s1);
            \path[->] (s2) edge[bend left=12] (s4);
            \path[->] (s2) edge[bend left=12] (s3);
            
            \path[->] (s3) edge[bend left=12] (s1);
            \path[->] (s3) edge[bend left=12] (s4);
            \path[->] (s3) edge[bend left=12] (s2);
            
            \path[->] (s4) edge[bend left=12] (s2);
            \path[->] (s4) edge[bend left=12] (s3);
            \path[->] (s4) edge[bend left=12] (s1);
        \end{tikzpicture}
        \caption{MDP}
    \end{subfigure}%
    \hfill
    \begin{subfigure}[b]{0.48\textwidth}
        \centering
        \begin{tikzpicture}[ 
            scale=1, auto, node distance=2cm, 
            every edge/.style={draw, <-, font=\small},
            state/.style={circle, fill=mylightgrey, draw=black, text=black}
            ]
        
            % Nodes
            \node[state] (s1) at (0,2) {$s_1$};
            \node[state] (s2) at (2,2) {$s_2$};
            \node[state] (s3) at (2,0) {$s_3$};
            \node[state] (s4) at (0,0) {$s_4$};
        
            \path[->] (s2) edge[loop above, out=25, in=65, looseness=8, draw=mylightgrey] (s2);
            \path[->] (s4) edge[loop below, out=205, in=245, looseness=8, draw=mylightgrey] (s4);
            
            \path[->] (s1) edge[bend left=12, draw=mylightgrey] (s2);
            \path[->] (s1) edge[bend left=12, draw=mylightgrey] (s3);
            \path[->] (s1) edge[bend left=12, draw=mylightgrey] (s4);
            
            \path[->] (s2) edge[bend left=12, draw=mylightgrey] (s1);
            \path[->] (s2) edge[bend left=12, draw=mylightgrey] (s4);
            \path[->] (s2) edge[bend left=12, draw=mylightgrey] (s3);
            
            \path[->] (s3) edge[bend left=12, draw=mylightgrey] (s1);
            
            \path[->] (s4) edge[bend left=12, draw=mylightgrey] (s2);
            \path[->] (s4) edge[bend left=12, draw=mylightgrey] (s3);
            \path[->] (s4) edge[bend left=12, draw=mylightgrey] (s1);

            \path[->] (s1) edge[loop above, out=115, in=155, looseness=8] (s1);
            \path[->] (s3) edge[loop below, out=295, in=335, looseness=8] (s3);
            \path[->] (s3) edge[bend left=12] (s4);
            \path[->] (s3) edge[bend left=12] (s2);
        \end{tikzpicture}
        \caption{Policy $\pol_1$}
    \end{subfigure}%

    \vspace{1cm}
    \vfill

    \begin{subfigure}[b]{0.48\textwidth}
        \centering
        \begin{tikzpicture}[ 
            scale=1, auto, node distance=2cm, 
            every edge/.style={draw, <-, font=\small},
            state/.style={circle, fill=mylightgrey, draw=black, text=black}
            ]
        
            % Nodes
            \node[state] (s1) at (0,2) {$s_1$};
            \node[state] (s2) at (2,2) {$s_2$};
            \node[state] (s3) at (2,0) {$s_3$};
            \node[state] (s4) at (0,0) {$s_4$};
            
            \path[->] (s3) edge[loop below, out=295, in=335, looseness=8, draw=mylightgrey] (s3);
            \path[->] (s4) edge[loop below, out=205, in=245, looseness=8, draw=mylightgrey] (s4);
            
            % Arcs between nodes
            \path[->] (s1) edge[bend left=12, draw=mylightgrey] (s2);
            \path[->] (s1) edge[bend left=12, draw=mylightgrey] (s3);
            
            \path[->] (s2) edge[bend left=12, draw=mylightgrey] (s1);
            \path[->] (s2) edge[bend left=12, draw=mylightgrey] (s4);
            
            \path[->] (s3) edge[bend left=12, draw=mylightgrey] (s1);
            \path[->] (s3) edge[bend left=12, draw=mylightgrey] (s4);
            \path[->] (s3) edge[bend left=12, draw=mylightgrey] (s2);
            
            \path[->] (s4) edge[bend left=12, draw=mylightgrey] (s2);
            \path[->] (s4) edge[bend left=12, draw=mylightgrey] (s3);
            \path[->] (s4) edge[bend left=12, draw=mylightgrey] (s1);
    
            \path[->] (s1) edge[loop above, out=115, in=155, looseness=8] (s1);
            \path[->] (s2) edge[loop above, out=25, in=65, looseness=8] (s2);
            \path[->] (s2) edge[bend left=12] (s3);
            \path[->] (s1) edge[bend left=12] (s4);    
        \end{tikzpicture}
        \caption{Policy $\pol_2$}
    \end{subfigure}%
    \hfill
    \begin{subfigure}[b]{0.48\textwidth}
        \centering
        \begin{tikzpicture}[ 
            scale=1, auto, node distance=2cm, 
            every edge/.style={draw, <-, font=\small},
            state/.style={circle, fill=mylightgrey, draw=black, text=black}
            ]
        
            % Nodes
            \node[state] (s1) at (0,2) {$s_1$};
            \node[state] (s2) at (2,2) {$s_2$};
            \node[state] (s3) at (2,0) {$s_3$};
            \node[state] (s4) at (0,0) {$s_4$};
            
            \path[->] (s2) edge[loop above, out=25, in=65, looseness=8, draw=mylightgrey] (s2);
            \path[->] (s3) edge[loop below, out=295, in=335, looseness=8, draw=mylightgrey] (s3);
            
            \path[->] (s1) edge[bend left=12, draw=mylightgrey] (s3);
            \path[->] (s1) edge[bend left=12, draw=mylightgrey] (s4);
            
            \path[->] (s2) edge[bend left=12, draw=mylightgrey] (s1);
            \path[->] (s2) edge[bend left=12, draw=mylightgrey] (s4);
            \path[->] (s2) edge[bend left=12, draw=mylightgrey] (s3);
            
            \path[->] (s3) edge[bend left=12, draw=mylightgrey] (s1);
            \path[->] (s3) edge[bend left=12, draw=mylightgrey] (s4);
            \path[->] (s3) edge[bend left=12, draw=mylightgrey] (s2);
            
            \path[->] (s4) edge[bend left=12, draw=mylightgrey] (s2);
            \path[->] (s4) edge[bend left=12, draw=mylightgrey] (s1);

            \path[->] (s1) edge[loop above, out=115, in=155, looseness=8] (s1);
            \path[->] (s1) edge[bend left=12] (s2);
            \path[->] (s4) edge[bend left=12] (s3);
            \path[->] (s4) edge[loop below, out=205, in=245, looseness=8] (s4);
        \end{tikzpicture}
        \caption{Policy $\pol_3$}
    \end{subfigure}%

    \vspace{1cm}
    
    \caption{The general Borkar-Meyn Theorem does not apply to MC-O-PI because convex combinations of the action values $q_{\pol_1}$, $q_{\pol_2}$ and $q_{\pol_3}$ can intersect the boundary between the regions of policies $\pol_1$, $\pol_2$ and $\pol_3$.
    This intersection introduces a suboptimal fixed point in the deterministic differential inclusion \eqref{eq: convex differential inclusion} that results from taking the convex hull, averaging out the noise and letting the learning rate tend to zero in the stochastic difference inclusion underlying MC-O-PI.}
    \label{fig: convex combi}
\end{figure}